\documentclass[11pt]{scrartcl}

\usepackage[stixtwo]{fontsetup}
\usepackage{amsmath}
\usepackage{tikz}

\begin{document}
\begin{enumerate}
	\item Bővítsük a törtünket a nevező konjugáltjával, majd végezzük el a
	      megfelelő egyszerűsítéseket:
	      \[
		      \frac{3 + 2i}{i - 4}
		      = \frac{(3 + 2i)(-4 - i)}{(-4 + i)(-4 - i)}
		      = \frac{-12 - 3i - 8i - 2i^2}{(-4)^2 + 1^2}
		      = -\frac{10}{17} - \frac{11}{17} i
		      \text.
	      \]

	\item \textbf{Szöveges gondolatmenet}

	      A $z = -8i$ egy olyan komplex szám, amelynek valós része
	      $\operatorname{Re} \{z\} = 0$, képzetes része pedig
	      $\operatorname{Im} \{z\} = -8 < 0$.

	      Tehát a hossza $r = |z| = 8$, argumentuma pedig
	      $\arg \{z\} = 270^\circ = 3\pi / 2$.

	      A harmadik gyökök hossza $\sqrt[3]{8} = 2$ lesz, az egyik gyök
	      argumentuma pedig pont az eredeti szám argumentumának harmada, vagyis
	      $\varphi_1 = 270^\circ / 3 = 90^\circ = \pi / 2$. A másik két
	      gyököt az előbb meghatározott gyök $120^\circ = 2\pi/3$, illetve
	      $240^\circ = 4\pi / 3$-mal való elforgatásával kaphatjuk meg. Tehát
	      a gyökök:
	      \[
		      \sqrt[3]{-8i} = \begin{cases}
			      2 \cdot e^{i \cdot \frac{\pi}{2}}
			      = 2i
			      \text,
			      \\
			      2 \cdot e^{i \cdot \left(\frac{\pi}{2} + \frac{2\pi}{3}\right)}
			      = 2 \cdot e^{i \cdot \frac{7\pi}{6}}
			      \text,
			      \\
			      2 \cdot e^{i \cdot \left(\frac{\pi}{2} + \frac{4\pi}{3}\right)}
			      = 2 \cdot e^{i \cdot \frac{11\pi}{6}}
			      \text.
		      \end{cases}
	      \]


	      \textbf{Másik megoldás}

	      $-8i = 8\left(\cos \frac{3\pi}{2} + i \sin \frac{3\pi}{2}\right)$, így
	      \[
		      w_k
		      = \sqrt[3]{8}\left(
		      \cos\left(\frac{3\pi/2 + 2k\pi}{3}\right)
		      + i \sin\left(\frac{3\pi/2 + 2k\pi}{3}\right)
		      \right),
		      \quad k \in \{0; 1; 2\}
		      \text.
	      \]
	      A gyökök tehát:
	      \[
		      \begin{aligned}
			      w_0 & = 2 \left(
			      \cos \frac{\pi}{2} + i \sin \frac{\pi}{2}
			      \right) = 2i
			      \text,
			      \\
			      w_1 & = 2 \left(
			      \cos \frac{7\pi}{6} + i \sin \frac{7\pi}{6}
			      \right)
			      \text,
			      \\
			      w_2 & = 2 \left(
			      \cos \frac{11\pi}{6} + i \sin \frac{11\pi}{6}
			      \right)
			      \text.
		      \end{aligned}
	      \]
	      A gyökök egy origó középpontú, $2$ sugarú körbe írt szabályos háromszög
	      csúcsai:
	      \begin{center}
		      \begin{tikzpicture}[scale=0.8, >=stealth]
			      \draw[->] (-2.8, 0) -- (2.8, 0) node[right] {$\operatorname{Re}$};
			      \draw[->] (0, -2.8) -- (0, 2.8) node[above] {$\operatorname{Im}$};
			      \draw[dashed] (0, 0) circle (2);
			      \node[above right] at (2, 0) {$2$};
			      \draw[thin] (90:2) -- (210:2) -- (330:2) -- cycle;
			      \fill (90:2) circle (2pt) node at (102:2.5) {$w_0$};
			      \fill (210:2) circle (2pt) node at (222:2.5) {$w_1$};
			      \fill (330:2) circle (2pt) node at (342:2.5) {$w_2$};
		      \end{tikzpicture}
	      \end{center}
\end{enumerate}
\end{document}
